% Agent documentation:
% Purpose: document survey-derived segment descriptions and implications.
% Downstream use: can support segment reasoning in the main text, but the small
% sample size limits evidential strength.
% Revision guardrails: preserve the caution around N=54 and weak cluster
% stability. Do not present these findings as representative without additional
% evidence.
\subsection{Spørreundersøkelse -- segmentbeskrivelser og implikasjoner}
\label{app:sporreundersokelse}

Basert på en klusteranalyse av spørreundersøkelsen (N\,=\,54) identifiseres
fire distinkte kundesegmenter. Forbeholdet er at utvalget er lavt for en
stabil klusterløsning, og at segment K2 (n\,=\,8) bør tolkes med særlig
varsomhet. Resultatene bør valideres kvalitativt, for eksempel gjennom
fokusgrupper.

\subsubsection*{Segmentbeskrivelser}

\textbf{K1 -- Miljøbevisste allroundere (11 stk, 20\,\%).}
Høyest miljøbevissthet og teknologiåpenhet av alle segmenter, kombinert med
høyt tidspress og prisbevissthet. Jevn kjønnsfordeling, blandet aldersprofil
med flest 18--25 og 46--55, og 55\,\% fulltidsansatte. 55\,\% har inntekt over
700\,k. Bor primært i Vestland og Oslo. 36\,\% vil bruke drone. Pålitelighet
(6,09) og lav pris (5,73) er viktigste kriterier, men miljø scorer også høyt
(5,00). Sikkerhet er den største bekymringen (73\,\%).

\textbf{K2 -- Lavt engasjerte skeptikere (8 stk, 15\,\%).}
Lavest på teknologiåpenhet og tidspress. Eldst profil -- 62\,\% er over 46 år,
50\,\% over 700\,k inntekt, 50\,\% bor i utkant av by eller landlig. 62\,\%
bestiller aldri hjemlevering. Kun 12\,\% vil bruke drone. Lavt prioritert
segment for Zipline -- vanskelig å konvertere og lite motivert for endring.

\textbf{K3 -- Travle pragmatikere (16 stk, 26\,\%).}
Høyest tidspress (5,07), middels teknologiåpenhet, lav miljøbevissthet. 64\,\%
er 18--25 år, primært studenter, bor i Trøndelag og Innlandet. Bestiller
relativt hyppig hjemlevering. Kun 29\,\% vil bruke drone. Rask levering (5,86)
og pålitelighet (6,00) er viktigst -- miljø er lite viktig (2,93).
Pålitelighet er den største bekymringen (57\,\%).

\textbf{K4 -- Unge prisbevisste tek-positive (21 stk, 39\,\%).}
Størst segment. Høy teknologiåpenhet og svært høy prissensitivitet. 90\,\% er
under 26 år, 71\,\% studenter, 52\,\% med inntekt under 100\,k. 81\,\% bor
sentralt i by, flest fra Trøndelag. 62\,\% vil bruke drone -- høyest av alle
segmenter. Lav pris (6,10) og pålitelighet (6,14) er viktigst, miljø lite
viktig (3,00). Pålitelighet er den største bekymringen (62\,\%).

\subsubsection*{Implikasjoner for Zipline}

\textbf{Primærmål (early adopters): K4 (39\,\%).}
Størst segment med høyest drone-vilje. Unge, urbane og teknologiåpne. Pris og
bekvemmelighet er nøkkelbudskapet, ikke miljø. Naturlig startsegment i
Trondheim.

\textbf{Sekundærmål: K1 (20\,\%).}
Kan nås med et kombinert budskap om teknologi, miljø og effektivitet. Høyere
inntekt gjør dem mindre prissensitive i praksis.

\textbf{Potensiell vekst: K3 (26\,\%).}
Tidspresset og aktive brukere av hjemlevering. Kan konverteres dersom
pålitelighet og hastighet kommuniseres tydelig.

\textbf{Deprioritert: K2 (15\,\%).}
Lav teknologiåpenhet og drone-vilje. Ikke lønnsomt å prioritere i en tidlig
markedsfase.
